\documentclass[11pt]{article}
\usepackage[a4paper,margin=25mm]{geometry}
\usepackage[T1]{fontenc}
\usepackage{lmodern}
\usepackage{booktabs}
\usepackage{amsmath}
\usepackage{graphicx}
\usepackage{hyperref}
\usepackage{xcolor}
\usepackage{listings}
\lstset{basicstyle=\ttfamily\small, columns=fullflexible, keepspaces=true,
  breaklines=true, frame=single, framerule=0.3pt, rulecolor=\color{gray}}
\newcommand{\code}[1]{\texttt{#1}}
\newcommand{\us}{\textmu s}

\title{\code{spe.py}: single-photoelectron spectra and the equalization factors}
\author{BGOTSCAN}
\date{2026-10-14}

\begin{document}
\maketitle

\section{Purpose}

The SiPMs (Hamamatsu S14160-6050HS, through a shaper into the digitizer)
produce dark pulses of one photoelectron, about 15 to 20 counts high and
24~ns wide at half maximum on 4 counts of noise per sample. Their height
is the gain of the channel in that run. It is needed to equalize the 64
channels before summing (\code{chanstat\_doc.tex} found four channels with
30 to 50\% more gain), to express the light yield in photoelectrons, and,
through its dependence on bias, to measure the breakdown voltage at each
temperature. There is no dedicated SPE run: the pulses are taken from the
pulse-free part of the scintillation records.

A threshold on single samples cannot separate these pulses from noise (a
first version of this script showed spectra cut by the threshold at every
bias). The search uses a matched filter built from the pulse shape
measured by \code{template.py}. The script walks one event in six steps,
as \code{selection.py}; its normal use is the \code{--summary} mode over
many events, which prints the per-channel single-photoelectron amplitudes
and the equalization factors. It is a Python script in \code{src/py/},
reading the raw HDF5 files with h5py; the shared definitions come from the
package \code{bgo} (\code{bgo\_doc.tex}).

\section{Usage}

\begin{lstlisting}
python3 src/py/spe.py RUN [options]
\end{lstlisting}

Interactive: after each step the script waits on the terminal; Enter goes
to the next step, \code{n} abandons the event and starts the next one,
\code{q} quits.

\begin{center}
\begin{tabular}{l l p{74mm}}
\toprule
option & default & meaning \\
\midrule
\code{--data DIR} & the data directory & as in \code{wfview.py} \\
\code{--events SPEC} & \code{0} & events to process: \code{a-b} or \code{a,b,c}, 0-based \\
\code{--channel C} & 0 & the channel shown in steps 1 to 3 and in the single-channel spectrum of step~6 \\
\code{--template FILE} & \code{src/spe\_template\_RUN.csv} & the pulse shape, as written by \code{template.py}; the default is the file of the run being processed, and the script stops if it is missing \\
\code{--region R} & \code{tail} & search region: \code{tail}, the last \code{--tail} \us{} of the record, or \code{all} \\
\code{--tail T} & 10 & length of the tail region in \us \\
\code{--nsigma K} & 4.5 & candidate threshold on the filtered waveform, in units of its noise \\
\code{--from-step N} & 1 & start the display at this step \\
\code{--save NAME} & off & no prompts; figures saved as \code{src/py/plots/RUN/NAME\_evE\_sN.png} \\
\code{--summary} & off & no prompts and no per-event figures; the step~6 text and figure after the last event, saved as \code{NAME\_spe.png} with \code{--save} \\
\code{--csv FILE} & off & write the per-channel table of step~6 to \code{FILE} \\
\code{--dump FILE} & off & write the candidates of every channel per event (time, $A$) and $\sigma_y$ to a JSON file (reference for the tests of the Julia pass) \\
\bottomrule
\end{tabular}
\end{center}

\section{Notation}

$w_{c,i}$ is the raw ADC value of channel $c$ at sample $i$ ($n_s =
3000$, $\Delta t = 8$~ns); $x_{c,i} = w_{c,i} - \operatorname{median}_i
w_{c,i}$, the channel minus its own baseline over the record. No
common-mode subtraction. $T_j$, $j = -H \dots H$ ($H = 8$), is the
template, normalised to $T_0 = 1$. $R$ is the search region.

\section{The steps}

\subsection{Step 1: read the channel}

Text: baseline, rms of $x$, the region in ns and samples, the template
file and its sum of squares. Figure: $x$ over the record with $R$ shaded.

\subsection{Step 2: matched filter and candidates}

The filtered waveform is the least-squares height of a pulse of shape $T$
at every sample,
\[
  y_i = \frac{\sum_j T_j\, x_{i+j}}{\sum_j T_j^2},
\]
in counts of the raw peak ($y_i$ equals the height of a pulse of shape $T$
centred at $i$). Its noise $\sigma_y$ is the clipped rms of $y$ over $R$
(samples beyond three times a first rms removed), per channel and event;
for white noise $\sigma_y = \sigma / \sqrt{\sum T_j^2}$, 2.6 counts for
$\sigma = 4$ and the template of \code{195K\_41V\_th20}. Candidates are the
local maxima of $y$ inside $R$ above $K \sigma_y$ (\code{--nsigma}, 4.5),
taken in decreasing order and at least $2H$ samples apart.

Text: $\sigma_y$, the threshold in counts, the number of candidates, their
times and heights. Figure: $x$ and $y$ over $R$, the threshold line, the
candidates marked.

\subsection{Step 3: the pulses}

For each candidate at $i_y$: its amplitude $A = y_{i_y}$, and the residual
rms of the fit, $\mathrm{rms}_j (x_{i_y + j} - A\,T_j)$, which says how
well the template describes it (about $\sigma$ for a true pulse, larger
for two overlapping ones). Text: one line per candidate: time, $A$,
residual. Figure: up to six candidates zoomed, the raw samples with
$A\,T$ drawn over them.

\subsection{Step 4: all 64 channels}

Steps 1 to 3 on every channel. Text: candidates per channel as an
$8\times8$ table, the total and the mean per channel. Figure: the
$8\times8$ grid of the filtered waveforms over $R$ with the candidates
marked.

\subsection{Step 5: rates}

$n_c / T_R$ per channel, with $T_R$ the length of $R$ in seconds. Text:
the mean rate and the channels more than three standard deviations from
it. Figure: the $8\times8$ map.

\subsection{Step 6: spectra, fit and equalization}

The amplitudes $A$ of all candidates are accumulated per channel over
the events processed so far. The expected spectrum of a channel: a peak
at the single-photoelectron amplitude $A_1$, a smaller one near $2 A_1$
from crosstalk, a tail from overlaps, and the low side cut by the
threshold $K \sigma_y$.

\paragraph{Fit.} A Gaussian is fitted to the histogram of $A$ (1-count
bins) between the threshold plus one $\sigma_y$ and 1.6 times the mode,
by least squares with Poisson weights, the mean bounded to the fit range
and the width to $[0.5\,\sigma_y, 1.6 \times \mathrm{mode}]$. Starting
the range at the threshold itself was tried, on the grounds that the cut
on $A$ only truncates the peak: the Gaussian then stretches over the
truncation and the noise spike of the first bins, fewer channels pass
the width bound (37 against 45 at 41~V) and the normal channels scatter
more; the range one $\sigma_y$ above the cut is kept. Its mean is $A_1$, its width
$w_1$; a fit is accepted if it has at least 30 pulses in the range, the
mean lies inside the range, and $w_1$ is between $0.5\,\sigma_y$ and
$A_1 / 2$. The fit is done for every channel and for all live channels
together (a live channel has at least one candidate; the dead ones of
\code{chanstat\_doc.tex} have none).

\paragraph{Equalization factor.} $g_c = A_1(c) / \operatorname{median}_c
A_1(c)$ over the channels with an accepted fit. Dividing channel $c$ by
$g_c$ puts all channels in the units of the median channel. The gain of
the run is the median $A_1$.

Text: for the selected channel and for all live channels, the number of
pulses, $A_1 \pm$ its error, $w_1$, the crosstalk ratio (pulses between
$1.5 A_1$ and $2.5 A_1$ over pulses below $1.5 A_1$); then one line per
channel: $n$, $A_1$, error, $w_1$, $g_c$, or ``no fit'' with the reason;
the median $A_1$ and the channels with $g_c$ more than 15\% from one.
Figure: four panels: the spectrum of the selected channel with its fit;
the spectrum of all live channels with its fit; the $8\times8$ map of
$A_1$; the $8\times8$ map of $g_c$. With \code{--csv}, the per-channel
table is written (channel, $n$, $A_1$, error, $w_1$, $g_c$).

\section{\code{spefit.py}: the fit on the output of the pass}

The pass of \code{JScan} (\code{jscan\_doc.tex}) applies steps 1 to 3 to
every channel of every event of a run and writes the candidates (event,
channel, time, $A$, residual) and $\sigma_y$ per channel and event to
\code{pass/RUN\_pass.h5}. \code{spefit.py} does step~6 on that table: the
spectra, the fits, the equalization factors. It reads only the pass
output, in a second, and uses the fit of \code{bgo.fit}, so that the peak
fit is defined once.

\begin{lstlisting}
python3 src/py/spefit.py RUN|FILE [--channel 48] [--min-pulses 30]
        [--dead LIST] [--save NAME] [--csv FILE]
\end{lstlisting}

\begin{center}
\begin{tabular}{l l p{76mm}}
\toprule
option & default & meaning \\
\midrule
\code{RUN|FILE} & & a run name, looked up as \code{pass/RUN\_pass.h5}, or a path \\
\code{--channel C} & 48 & the channel whose spectrum is shown alone \\
\code{--min-pulses N} & 30 & channels with fewer candidates are not fitted \\
\code{--dead LIST} & the twelve of \code{chanstat\_doc.tex} & excluded from the reference and the maps \\
\code{--save NAME} & off & figures saved as \code{src/py/plots/RUN/NAME\_spefit.png} and \code{NAME\_spefit\_channels.png}; otherwise shown \\
\code{--csv FILE} & off & write the per-channel table (channel, $n$, $A_1$, error, $w_1$, $g_c$, rate) \\
\bottomrule
\end{tabular}
\end{center}

Per channel, the threshold used by the pass is $K$ times the mean of
$\sigma_y$ over the events ($K$ from the file metadata), and the fit is
that of step~6: a bounded Gaussian between the threshold plus one
$\sigma_y$ and 1.6 times the mode. The reference for $g_c$ is the median
$A_1$ over the live channels with an accepted fit, dead channels
excluded. The dark rate of a channel is its number of candidates divided
by the total search time (events times the length of the search region).

Text: the metadata of the pass (run, template, threshold, events), the
fit of the selected channel and of all live channels together, one line
per channel, the median $A_1$, the channels with $g_c$ more than 15\% from
one. Figures: the four panels of step~6, and an $8\times8$ grid of the
per-channel spectra with their fits, to judge each fit by eye.

\section{Modes}

Interactive, one prompt per step; \code{--from-step 6}, one prompt per
event showing the growing spectra; \code{--save}, no prompts, every
figure saved; \code{--summary}, nothing per event, step~6 after the last
event. The summary over a full run of 5000 events is the normal use.

\section{Examples}

\begin{lstlisting}
# walk through event 0, channel 48, of the 41 V run
python3 src/py/spe.py 195K_41V_th20 --channel 48

# the equalization table of the 41 V run, full statistics
python3 src/py/spe.py 195K_41V_th20 --events 0-4999 --summary --save spe195 \
        --csv src/spe_195K_41V_th20.csv

# the same for a run without its own template, using the 41 V one
python3 src/py/spe.py 196K_40V_th20 --events 0-4999 --summary \
        --template src/spe_template_195K_41V_th20.csv
\end{lstlisting}

\end{document}
